\chapter{Pressure, Free Energy, and Energy Fluctuations}

\lectureoverview{We derive pressure from the partition function without
assuming a kinetic picture of molecules striking a wall.  A geometric lemma
converts an adiabatic derivative into fixed-temperature derivatives, the
Helmholtz free energy supplies the bridge to \(\log Z\), and the ideal-gas law
follows as the first application.  We then use the same partition function to
derive the equilibrium energy fluctuation and its relation to heat capacity.}

\section{Let the Mathematics Test the Intuition}

Thermodynamics has a peculiar reputation.  Its calculus is elementary, yet
the choice of what to differentiate and what to hold fixed can be remarkably
subtle.  Our problem is deliberately narrow: calculate the pressure of an
ideal gas.  The route to that answer, however, is intended to be general
enough for systems in which the familiar molecular picture is no longer
reliable.

The kinetic story is easy to tell.  Molecules strike a wall and reverse the
normal component of their momentum.  Momentum delivered per unit time gives
the force, and force per unit area gives the pressure.  We already know that
the mean kinetic energy per molecule in three dimensions is
\begin{equation}
 \left\langle \frac{p^2}{2m}\right\rangle=\frac{3}{2}T
 \qquad (k_B=1).
 \label{eq:sm5-mean-kinetic-energy}
\end{equation}
With an isotropic velocity distribution one can estimate the collision rate,
the momentum transfer in each collision, and the resulting pressure.

But the wall itself seems to spoil isotropy.  Why should molecules very near
the boundary have the same distribution of speeds and directions as
molecules deep in the interior?  In equilibrium they do, to the accuracy
relevant here, but that conclusion should follow from statistical mechanics
rather than be smuggled into the calculation as intuition.  Collisions among
the molecules make the kinetic argument still less transparent.  The
noninteracting gas is easy precisely because it is an unusually simple
system.

There are occasions in physics when the safest procedure is to suppress the
picture temporarily and ride the mathematics.  One identifies the quantity
to be found, manipulates the general relations, and waits until a recognizable
formula appears.  Intuition returns at the end, but then the mathematics is
testing the intuition rather than the other way around.  The pleasure of
thermodynamics lies in the surprising relations that emerge from this
procedure.

The desired relation is an equation of state,
\begin{equation}
 P=P(T,V,N),
 \label{eq:sm5-equation-of-state-target}
\end{equation}
where \(N\) is the amount of material, \(T\) the temperature, and \(V\) the
volume.  A satisfactory derivation should continue to make sense when the
particles interact, or when the material is a liquid, plasma, or solid.  Only
the final evaluation of the partition function will specialize to an ideal
gas.

\section{Entropy and the Helmholtz Free Energy}

Begin by retrieving the canonical formulas that will be needed later.  For
microstates labelled by \(i\),
\begin{equation}
 \begin{aligned}
 S&=-\sum_i p_i\log p_i,\\
 p_i&=\frac{e^{-\beta E_i}}{Z},\\
 Z&=\sum_i e^{-\beta E_i}.
 \end{aligned}
 \label{eq:sm5-canonical-definitions}
\end{equation}
The logarithm of the probability is
\begin{equation}
 \log p_i=-\beta E_i-\log Z.
 \label{eq:sm5-log-probability}
\end{equation}
Substitution into the entropy gives
\begin{align}
 S
 &=-\sum_i p_i\log p_i \notag\\
 &=\beta\sum_i p_iE_i+\log Z\sum_i p_i \notag\\
 &=\beta E+\log Z.
 \label{eq:sm5-entropy-log-z}
\end{align}
Here \(E=\sum_i p_iE_i\) is the mean energy.

\begin{classroomqa}
\audiencequestion{Why does the second term give only \(\log Z\), rather
than the number of states times \(\log Z\)?}
\lecturerresponse{The coefficient being summed is \(p_i\), not \(1\).
Normalization gives \(\sum_i p_i=1\).  Equivalently, the numerator
\(\sum_i e^{-\beta E_i}=Z\) cancels the \(Z\) in the denominator.  The
number of states never appears in this step.
\lecturetimestamp{00:14:25--00:14:55}}
\end{classroomqa}

In the energy units used through most of the lecture,
\(\beta=1/T\).  Multiplying \eqref{eq:sm5-entropy-log-z} by \(T\) and
rearranging gives
\begin{equation}
 E-TS=-T\log Z.
 \label{eq:sm5-helmholtz-identity}
\end{equation}
The combination on the left occurs so often that it deserves a name:
\begin{equation}
 A\equiv E-TS,
 \qquad
 A=-T\log Z.
 \label{eq:sm5-helmholtz-definition}
\end{equation}
It is the Helmholtz free energy.  The name is less important at this moment
than the identity itself; we shall need the identity when pressure is
converted into a derivative of the partition function.

\begin{figure}[t]
\centering
\includegraphics[width=0.90\textwidth]{lecture_05_figure_03.png}
\caption{The canonical entropy relation reorganized into
\(A=E-TS=-T\log Z\).}
\lectureframe{00:16:50}
\end{figure}

\section{Control Parameters and Their Responses}

A control parameter is a macroscopic quantity that an experimenter can
change from outside.  Moving a piston changes the volume without manipulating
individual molecules.  Applied electric and magnetic fields are other
examples.  In a more general ensemble, the number of particles can also be
allowed to change.

Thermodynamic variables naturally come in conjugate pairs.  Volume is the
control parameter in the present problem, and pressure is its response.  We
change \(V\); the pressure tells us how the system energetically resists that
change.  Later the derivative of the energy will make this pairing precise.

\begin{classroomqa}
\audiencequestion{Can pressure itself be regarded as the control parameter?}
\lecturerresponse{One may design an experiment that fixes the external
pressure, but in the construction used here the piston position fixes the
volume and the internal pressure is the response.  Pressure and volume are
conjugate variables; which member is controlled depends on the ensemble and
apparatus.
\lecturetimestamp{00:19:22--00:20:01}}
\end{classroomqa}

\begin{classroomqa}
\audiencequestion{Are all the molecules assumed to be identical, or is the
argument restricted to a gas?}
\lecturerresponse{No such restriction is needed for the thermodynamic
argument.  The cylinder may contain several species, a liquid, or even a
solid being compressed.  The ideal-gas partition function will be a special
application only after the general pressure formula has been established.
\lecturetimestamp{00:20:02--00:20:39}}
\end{classroomqa}

\begin{classroomqa}
\audiencequestion{Could the amount of material be changed as another
external operation?}
\lecturerresponse{Yes.  Particle number is another possible control
parameter.  Treating it dynamically leads to its own conjugate variable, the
chemical potential.  For the calculation tonight, \(N\) remains fixed.
\lecturetimestamp{00:20:39--00:20:57}}
\end{classroomqa}

\begin{classroomqa}
\audiencequestion{Can liquid water and steam be present simultaneously in
equilibrium?}
\lecturerresponse{At the coexistence temperature and pressure they can.
Gravity may place the liquid below the vapor, yet the two phases can still
share one temperature and remain in thermal equilibrium.  The same general
statistical-mechanical framework applies.
\lecturetimestamp{00:20:59--00:21:36}}
\end{classroomqa}

For the next argument, regard energy and entropy as dependent functions of
two independently adjustable variables:
\begin{equation}
 E=E(T,V),
 \qquad
 S=S(T,V).
 \label{eq:sm5-dependent-functions}
\end{equation}
Temperature and volume provide coordinates on a thermodynamic state space.
Energy and entropy assign numbers to every equilibrium point in that space.

\section{The Constant-Entropy Derivative}

The derivative we ultimately need is unusual:
\[
 \left(\frac{\partial E}{\partial V}\right)_S.
\]
Although \(T\) and \(V\) are the independent coordinates, this derivative
asks us to change \(V\) while constraining the dependent function \(S(T,V)\)
to remain constant.  It is a derivative along a curve in the
\((T,V)\)-plane, not along a coordinate axis.

\begin{classroomqa}
\audiencequestion{How can \(T\) and \(V\) be independent if a derivative
\(dT/dV\) is later used?}
\lecturerresponse{They are independent on the full state space.  Once we
impose the constraint \(S(T,V)=\mathrm{constant}\), the allowed states lie
on a one-dimensional contour.  Along that contour, \(T\) becomes a function
of \(V\), and its slope \(dT/dV\) is well defined.
\lecturetimestamp{01:01:18--01:02:09}}
\end{classroomqa}

\begin{lemma}
If \(E(T,V)\) and \(S(T,V)\) are differentiable, then wherever
\((\partial S/\partial T)_V\neq0\),
\begin{equation}
 \left(\frac{\partial E}{\partial V}\right)_S
 =
 \left(\frac{\partial E}{\partial V}\right)_T
 -
 \left(\frac{\partial E}{\partial S}\right)_V
 \left(\frac{\partial S}{\partial V}\right)_T.
 \label{eq:sm5-calculus-lemma}
\end{equation}
\end{lemma}

\begin{proof}
For an arbitrary infinitesimal displacement,
\begin{equation}
 dE=
 \left(\frac{\partial E}{\partial V}\right)_T dV
 +
 \left(\frac{\partial E}{\partial T}\right)_V dT.
 \label{eq:sm5-total-differential-e}
\end{equation}
At fixed \(V\), the chain rule permits \(S\) to be used as an intermediate
variable:
\begin{equation}
 \left(\frac{\partial E}{\partial T}\right)_V
 =
 \left(\frac{\partial E}{\partial S}\right)_V
 \left(\frac{\partial S}{\partial T}\right)_V.
 \label{eq:sm5-chain-rule-e-s-t}
\end{equation}

Now restrict the displacement to a contour on which \(S\) is constant.
Along that contour,
\begin{equation}
 0=dS=
 \left(\frac{\partial S}{\partial V}\right)_T dV
 +
 \left(\frac{\partial S}{\partial T}\right)_V dT.
 \label{eq:sm5-constant-entropy-differential}
\end{equation}
Solving for the contour slope gives
\begin{equation}
 \left(\frac{dT}{dV}\right)_S
 =
 -
 \frac{\left(\frac{\partial S}{\partial V}\right)_T}
 {\left(\frac{\partial S}{\partial T}\right)_V}.
 \label{eq:sm5-entropy-contour-slope}
\end{equation}
Divide \eqref{eq:sm5-total-differential-e} by \(dV\), use
\eqref{eq:sm5-chain-rule-e-s-t}, and insert
\eqref{eq:sm5-entropy-contour-slope}:
\begin{align}
 \left(\frac{\partial E}{\partial V}\right)_S
 &=
 \left(\frac{\partial E}{\partial V}\right)_T
 +
 \left(\frac{\partial E}{\partial S}\right)_V
 \left(\frac{\partial S}{\partial T}\right)_V
 \left(\frac{dT}{dV}\right)_S \notag\\
 &=
 \left(\frac{\partial E}{\partial V}\right)_T
 -
 \left(\frac{\partial E}{\partial S}\right)_V
 \left(\frac{\partial S}{\partial V}\right)_T.
\end{align}
The factor \((\partial S/\partial T)_V\) has cancelled, proving the result.
\end{proof}

\begin{figure}[t]
\centering
\includegraphics[width=0.90\textwidth]{lecture_05_figure_04.png}
\caption{The constant-entropy contour calculation at the cancellation that
completes the derivative lemma.}
\lectureframe{00:36:10}
\end{figure}

The geometric picture is not ornamental.  It explains why an apparently
strange fixed-\(S\) derivative can be evaluated with ordinary partial
derivatives in the independent coordinates \(T\) and \(V\).  Writing
\(\delta E/\delta V\) during the construction is a useful reminder that the
symbols represent ratios of small changes along a specified path.

\begin{classroomqa}
\audiencequestion{Is the expansion of \(dE\) just the ordinary definition
of partial derivatives?}
\lecturerresponse{Yes.  Nothing special occurs in the total differential.
The only extra ingredient is the later restriction to the contour
\(dS=0\), which determines the relation between \(dT\) and \(dV\).
\lecturetimestamp{00:29:26--00:29:32}}
\end{classroomqa}

\begin{classroomqa}
\audiencequestion{Does the cancellation require
\((\partial S/\partial T)_V\) to be nonzero?}
\lecturerresponse{Yes.  For an ordinary stable system,
\[
 \left(\frac{\partial S}{\partial T}\right)_V
 =\frac{C_V}{T}>0,
\]
so the division is legitimate.  Exceptional points require separate care.
At fixed Hamiltonian, raising the temperature broadens the canonical energy
distribution and normally increases the entropy; lowering the temperature
narrows it.  This is the physically correct orientation of the monotonicity
argument.
\lecturetimestamp{00:36:50--00:37:47}}
\end{classroomqa}

\section{Piston Work and the Adiabatic Condition}

We now give the derivative a physical meaning.  Consider a piston with
cross-sectional area \(\mathcal A\).  If it moves outward by \(dx\), then
\begin{equation}
 dV=\mathcal A\,dx.
 \label{eq:sm5-piston-volume}
\end{equation}
The gas exerts force \(P\mathcal A\) on the piston, so the work done by the
gas is
\begin{equation}
 \delta W_{\mathrm{by}}=P\mathcal A\,dx=P\,dV.
 \label{eq:sm5-piston-work}
\end{equation}
If no heat enters the gas, energy conservation gives
\begin{equation}
 dE=-P\,dV.
 \label{eq:sm5-adiabatic-work}
\end{equation}
The sign is physical: during expansion the gas raises or accelerates the
piston, and its own energy decreases.

\begin{figure}[t]
\centering
\includegraphics[width=0.90\textwidth]{lecture_05_figure_05.png}
\caption{The piston geometry and the completed work relation
\(dE=-P\,dV\).}
\lectureframe{00:43:10}
\end{figure}

Two qualifications are essential.  The piston must move slowly enough for
the system to remain arbitrarily close to equilibrium, and the walls must be
insulated so that energy is not supplied as heat.  A rapid pull can be so
fast that no molecule reaches the receding piston during the motion.  A
rapid compression launches shock waves.  Neither process is represented by
a sequence of equilibrium states.

The word \emph{adiabatic} is sometimes used for no heat transfer and
sometimes for slow change.  In the reversible process required here, both
ideas enter: the change is quasistatic and thermally insulated.  Under those
conditions it is also isentropic.

\begin{classroomqa}
\audiencequestion{Why does an adiabatic change keep the entropy fixed?}
\lecturerresponse{View the system quantum mechanically.  Its many-body
energy levels depend on the volume.  When \(V\) changes slowly, the adiabatic
theorem carries each occupied level continuously into its successor without
transitions between level labels.  The energies \(E_i(V)\) move, but the
occupation probabilities \(p_i\) remain fixed.  Since
\(S=-\sum_i p_i\log p_i\) depends on those probabilities, the entropy remains
constant.  Heat exchange or irreversible motion would invalidate this
conclusion.
\lecturetimestamp{00:43:37--00:49:15}}
\end{classroomqa}

\begin{figure}[t]
\centering
\includegraphics[width=0.90\textwidth]{lecture_05_figure_06.png}
\caption{Many-body energy levels moving with the volume while their
occupation probabilities retain their level labels.}
\lectureframe{00:48:50}
\end{figure}

Equation \eqref{eq:sm5-adiabatic-work}, evaluated along this
constant-entropy path, identifies pressure as
\begin{equation}
 P=-\left(\frac{\partial E}{\partial V}\right)_S.
 \label{eq:sm5-pressure-fixed-entropy}
\end{equation}
This is the conjugate relation between pressure and volume.

\begin{classroomqa}
\audiencequestion{What happens if the piston is moved rapidly instead?}
\lecturerresponse{The system leaves equilibrium and its detailed dynamics
matter.  Sudden compression generates a shock.  In an ideal free expansion,
the piston can recede before any molecule strikes it, so no work is done and
the energy remains fixed even though the volume changes.  The gas later
relaxes to equilibrium with increased entropy.  Such processes cannot be
described by the single equilibrium derivative in
\eqref{eq:sm5-pressure-fixed-entropy}.
\lecturetimestamp{00:58:30--01:01:15}}
\end{classroomqa}

\begin{classroomqa}
\audiencequestion{What quantitative condition decides whether a change is
slow?}
\lecturerresponse{The external variation time must be long compared with
the relevant internal relaxation times.  For a gas, many molecular
collisions and many wall impacts should occur during the macroscopic change.
The precise time scale depends on the material and geometry, while the
quasistatic limit itself is general.
\lecturetimestamp{01:02:09--01:03:20}}
\end{classroomqa}

\section{From Adiabatic Energy to the Partition Function}

The physical formula \eqref{eq:sm5-pressure-fixed-entropy} is not yet the
easy one to calculate.  Canonical statistical mechanics naturally supplies
functions of temperature, because \(T\) appears directly in the Boltzmann
weight.  The calculus lemma converts the awkward fixed-entropy derivative
into that calculational language:
\begin{align}
 P
 &=
 -\left(\frac{\partial E}{\partial V}\right)_T
 +
 \left(\frac{\partial E}{\partial S}\right)_V
 \left(\frac{\partial S}{\partial V}\right)_T.
 \label{eq:sm5-pressure-lemma-substitution}
\end{align}
At fixed volume the first law gives \(dE=T\,dS\), and therefore
\begin{equation}
 \left(\frac{\partial E}{\partial S}\right)_V=T.
 \label{eq:sm5-e-s-temperature}
\end{equation}
Thus
\begin{equation}
 P=
 -\left(\frac{\partial E}{\partial V}\right)_T
 +
 T\left(\frac{\partial S}{\partial V}\right)_T.
 \label{eq:sm5-pressure-fixed-temperature-pieces}
\end{equation}

\begin{classroomqa}
\audiencequestion{Where does the minus sign in the pressure formula come
from?}
\lecturerresponse{It originates in work done by the gas:
\(dE=-P\,dV\) during adiabatic expansion.  The calculus lemma itself has a
minus sign between its two terms; the overall minus in
\(P=-(\partial E/\partial V)_S\) reverses both of them, producing
\eqref{eq:sm5-pressure-fixed-temperature-pieces}.
\lecturetimestamp{00:52:40--00:53:14}}
\end{classroomqa}

Because \(T\) is held fixed in both derivatives,
\begin{align}
 P
 &=
 -\left(\frac{\partial}{\partial V}(E-TS)\right)_T \notag\\
 &=
 -\left(\frac{\partial A}{\partial V}\right)_T.
 \label{eq:sm5-pressure-free-energy}
\end{align}
Finally insert \(A=-T\log Z\):
\begin{equation}
 \boxed{
 P=
 T\left(\frac{\partial\log Z}{\partial V}\right)_T
 }.
 \label{eq:sm5-pressure-log-z}
\end{equation}

\begin{figure}[t]
\centering
\includegraphics[width=0.90\textwidth]{lecture_05_figure_07.png}
\caption{The fixed-temperature compression of the argument:
\(P=-(\partial A/\partial V)_T
=T(\partial\log Z/\partial V)_T\).}
\lectureframe{00:55:55}
\end{figure}

Nothing in \eqref{eq:sm5-pressure-log-z} assumes an ideal gas.  If the
partition function can be calculated for an interacting fluid, plasma, or
solid, the same derivative gives its pressure.  More generally, every
external control parameter has a conjugate response obtained by
differentiating an appropriate thermodynamic potential.  The partition
function packages many such responses: differentiating \(\log Z\) with
respect to \(\beta\) gives energy, while differentiating it with respect to
volume gives pressure.

\section{The Ideal-Gas Equation of State}

Only now specialize to \(N\) noninteracting particles.  With the classical
phase-space normalization retained,
\begin{equation}
 Z_N=
 \frac{1}{N!h^{3N}}
 \int d^{3N}q\,d^{3N}p\,
 \exp\!\left[
 -\beta\sum_{a=1}^{N}\frac{\mathbf p_a^2}{2m}
 \right].
 \label{eq:sm5-ideal-gas-partition-integral}
\end{equation}
There are \(3N\) position coordinates and \(3N\) momentum coordinates.  The
Hamiltonian is independent of position, so every particle contributes one
factor of the container volume:
\begin{equation}
 Z_N=
 \frac{V^N}{N!h^{3N}}I_p(\beta),
 \label{eq:sm5-ideal-gas-factorization}
\end{equation}
where the entire momentum integral \(I_p(\beta)\) is independent of \(V\).

\begin{classroomqa}
\audiencequestion{Why are there \(3N\) position integrals and \(3N\)
momentum integrals?}
\lecturerresponse{Each of the \(N\) particles has three spatial coordinates
and three conjugate momentum components.  The count is \(3N\) in each half
of phase space, not \(N\) particles plus another \(3N\) particles.
\lecturetimestamp{01:04:19--01:04:28}}
\end{classroomqa}

Taking a logarithm separates the volume dependence:
\begin{equation}
 \log Z_N
 =
 N\log V-\log N!-3N\log h+\log I_p(\beta).
 \label{eq:sm5-log-ideal-gas-partition}
\end{equation}
The factors \(1/N!\), \(h^{-3N}\), and \(I_p(\beta)\) are all independent of
volume.  They matter for absolute entropy and other questions, but not for
this particular derivative.  Hence
\begin{equation}
 \left(\frac{\partial\log Z_N}{\partial V}\right)_T
 =\frac{N}{V}.
 \label{eq:sm5-log-z-volume-derivative}
\end{equation}
Equation \eqref{eq:sm5-pressure-log-z} immediately gives
\begin{equation}
 P=\frac{NT}{V},
 \qquad
 \boxed{PV=NT}.
 \label{eq:sm5-ideal-gas-law}
\end{equation}

\begin{classroomqa}
\audiencequestion{Does \(N\) still count the replicas used to derive the
Boltzmann distribution?}
\lecturerresponse{No.  The replica construction was a counting device in
the earlier derivation and is no longer present.  Here \(N\) is the number of
physical particles in this one gas sample.
\lecturetimestamp{01:07:38--01:08:05}}
\end{classroomqa}

Defining the number density
\begin{equation}
 \rho=\frac{N}{V}
 \label{eq:sm5-number-density}
\end{equation}
gives the equivalent form
\begin{equation}
 P=\rho T.
 \label{eq:sm5-density-equation-of-state}
\end{equation}

\begin{classroomqa}
\audiencequestion{Should the density form be obtained by dividing by \(N\)
or by \(V\)?}
\lecturerresponse{Divide \(PV=NT\) by \(V\):
\(P=(N/V)T=\rho T\).  Here \(\rho\) is number density, not mass density.
\lecturetimestamp{01:08:45--01:09:16}}
\end{classroomqa}

\begin{figure}[t]
\centering
\includegraphics[width=0.90\textwidth]{lecture_05_figure_08.png}
\caption{The volume derivative of \(\log Z\) and the resulting ideal-gas
equation of state \(PV=NT\), or \(P=\rho T\).}
\lectureframe{01:09:25}
\end{figure}

Restoring laboratory temperature \(T_K\),
\begin{equation}
 PV=Nk_BT_K.
 \label{eq:sm5-ideal-gas-law-laboratory}
\end{equation}
The familiar result agrees with kinetic intuition, but its derivation now
rests on a framework that survives far beyond the free gas.

\section{A Second Dividend: Equilibrium Fluctuations}

The pressure calculation is complete, but one promised question remains:
what else can be extracted from the partition function?  Consider a system
in thermal contact with a heat bath.  Its energy is not fixed.  Energy enters
and leaves, so repeated measurements fluctuate around the canonical mean.

\begin{classroomqa}
\audiencequestion{Can we now calculate the fluctuations that were promised?}
\lecturerresponse{Yes.  The extra time allows a second demonstration of the
same principle: derivatives of \(Z\) encode measurable thermodynamic
information.
\lecturetimestamp{01:10:37--01:10:50}}
\end{classroomqa}

\subsection{Variance and Standard Deviation}

For a random variable \(x\) with \(\langle x\rangle=0\), a measure of the
distribution's width is \(\langle x^2\rangle\).  For a general distribution,
shift the variable by its mean:
\begin{equation}
 \delta x=x-\langle x\rangle.
 \label{eq:sm5-centered-variable}
\end{equation}
Its mean vanishes,
\begin{equation}
 \langle\delta x\rangle
 =
 \langle x\rangle-\langle x\rangle=0,
 \label{eq:sm5-centered-mean}
\end{equation}
and its mean square is
\begin{align}
 (\Delta x)^2
 &\equiv
 \left\langle\left(x-\langle x\rangle\right)^2\right\rangle \notag\\
 &=
 \langle x^2\rangle
 -2\langle x\rangle^2
 +\langle x\rangle^2 \notag\\
 &=
 \langle x^2\rangle-\langle x\rangle^2.
 \label{eq:sm5-variance}
\end{align}
This is the variance.  Its positive square root \(\Delta x\) is the standard
deviation.

\begin{classroomqa}
\audiencequestion{Why is the average of
\(x-\langle x\rangle\) exactly zero?}
\lecturerresponse{Averages are linear, and \(\langle x\rangle\) is already
a number.  Therefore
\(\langle x-\langle x\rangle\rangle
=\langle x\rangle-\langle x\rangle=0\).
\lecturetimestamp{01:14:02--01:14:22}}
\end{classroomqa}

\begin{classroomqa}
\audiencequestion{What is the distinction among squared fluctuation,
variance, root-mean-square width, and standard deviation?}
\lecturerresponse{The variance is
\((\Delta x)^2=\langle x^2\rangle-\langle x\rangle^2\).
The standard deviation is its square root, \(\Delta x\).  For a zero-mean
variable the standard deviation is also its root-mean-square value.
\lecturetimestamp{01:16:39--01:18:38}}
\end{classroomqa}

\begin{figure}[t]
\centering
\includegraphics[width=0.90\textwidth]{lecture_05_figure_09.png}
\caption{The shifted variable and the expansion that produces
\((\Delta x)^2=\langle x^2\rangle-\langle x\rangle^2\).}
\lectureframe{01:17:20}
\end{figure}

\subsection{Energy Moments from the Partition Function}

For the canonical ensemble,
\begin{align}
 \langle E\rangle
 &=
 \frac{1}{Z}\sum_i E_i e^{-\beta E_i} \notag\\
 &=
 -\frac{1}{Z}\frac{\partial Z}{\partial\beta}
 =
 -\frac{\partial\log Z}{\partial\beta}.
 \label{eq:sm5-mean-energy-from-z}
\end{align}
Multiplication by one factor of \(E_i\) is produced by one negative
\(\beta\)-derivative.  A second derivative generates \(E_i^2\):
\begin{equation}
 \langle E^2\rangle
 =
 \frac{1}{Z}\sum_i E_i^2e^{-\beta E_i}
 =
 \frac{1}{Z}\frac{\partial^2 Z}{\partial\beta^2}.
 \label{eq:sm5-second-energy-moment}
\end{equation}
The energy variance is consequently
\begin{align}
 (\Delta E)^2
 &=
 \langle E^2\rangle-\langle E\rangle^2 \notag\\
 &=
 \frac{1}{Z}\frac{\partial^2Z}{\partial\beta^2}
 -
 \frac{1}{Z^2}
 \left(\frac{\partial Z}{\partial\beta}\right)^2 \notag\\
 &=
 \frac{\partial^2\log Z}{\partial\beta^2}
 =
 -\frac{\partial\langle E\rangle}{\partial\beta}.
 \label{eq:sm5-energy-variance-beta}
\end{align}
Once again, a complicated-looking observable becomes a derivative of
\(\log Z\).

\begin{figure}[t]
\centering
\includegraphics[width=0.90\textwidth]{lecture_05_figure_10.png}
\caption{The second \(\beta\)-derivative of \(\log Z\) identified with the
energy variance.}
\lectureframe{01:29:00}
\end{figure}

\subsection{Heat Capacity and the Size of Fluctuations}

In lecture units \(T=1/\beta\), so
\begin{equation}
 -\frac{\partial}{\partial\beta}
 =
 T^2\frac{\partial}{\partial T}.
 \label{eq:sm5-beta-temperature-derivative}
\end{equation}
At fixed volume define the heat capacity of the whole sample by
\begin{equation}
 C_V\equiv
 \left(\frac{\partial\langle E\rangle}{\partial T}\right)_V.
 \label{eq:sm5-heat-capacity}
\end{equation}
Equation \eqref{eq:sm5-energy-variance-beta} becomes
\begin{equation}
 \boxed{(\Delta E)^2=T^2C_V}
 \qquad (k_B=1).
 \label{eq:sm5-energy-fluctuation-lecture-units}
\end{equation}
For laboratory temperature \(T_K\), where
\(\beta=(k_BT_K)^{-1}\),
\begin{equation}
 \boxed{(\Delta E)^2=k_BT_K^2C_V}.
 \label{eq:sm5-energy-fluctuation-laboratory}
\end{equation}

\begin{figure}[t]
\centering
\includegraphics[width=0.90\textwidth]{lecture_05_figure_11.png}
\caption{The conversion from the \(\beta\)-derivative to the fixed-volume
heat capacity, including the restored factor of \(k_B\).}
\lectureframe{01:31:10}
\end{figure}

\begin{classroomqa}
\audiencequestion{Must the derivative specify fixed volume rather than fixed
pressure?}
\lecturerresponse{Yes.  This canonical derivation keeps the Hamiltonian's
external control parameters fixed, including \(V\).  The resulting quantity
is \(C_V\).  A constant-pressure ensemble allows volume and work to
fluctuate and requires the corresponding thermodynamic potential; one
cannot simply replace \(C_V\) by \(C_P\) in
\eqref{eq:sm5-energy-fluctuation-laboratory}.
\lecturetimestamp{01:29:01--01:29:28; 01:32:50--01:33:21}}
\end{classroomqa}

\begin{classroomqa}
\audiencequestion{Is \(C_V\) a material constant?}
\lecturerresponse{Not in general.  It may depend on temperature, volume,
phase, and composition.  Moreover, \(C_V\) here is the heat capacity of the
entire sample and is extensive.  A specific heat is obtained only after
dividing by mass, particle number, or amount of substance.
\lecturetimestamp{01:31:47--01:32:21; 01:33:23--01:34:55}}
\end{classroomqa}

For ordinary extensive matter, \(C_V\) is proportional to \(N\), so
\((\Delta E)^2\) is proportional to \(N\) and \(\Delta E\) to
\(\sqrt{N}\).  The mean energy itself is proportional to \(N\).  Therefore
\begin{equation}
 \frac{\Delta E}{\langle E\rangle}\sim\frac{1}{\sqrt{N}}.
 \label{eq:sm5-relative-fluctuation-scaling}
\end{equation}
Macroscopic relative fluctuations are tiny, while microscopic samples
display them directly.  The appearance of \(k_B\) in
\eqref{eq:sm5-energy-fluctuation-laboratory} marks this connection between
thermodynamic observables and molecular discreteness; equivalently, it
connects macroscopic units to Avogadro-scale counting.

Pressure and fluctuations have now emerged from the same object:
\[
 P=T\left(\frac{\partial\log Z}{\partial V}\right)_T,
 \qquad
 (\Delta E)^2=\frac{\partial^2\log Z}{\partial\beta^2}.
\]
This is the central lesson.  Once the correct constraints are identified,
the partition function does not merely normalize probabilities.  Its
derivatives organize the measurable thermodynamics of the system.
